%==============================================================================
% Curriculum Vitae — Juan José Jiménez Guardo
% Estilo: Harvard / ATS-friendly (sin fotografía, una sola columna)
% Compilador requerido: XeLaTeX o LuaLaTeX (por el uso de fontspec)
%
% ESTRUCTURA MODULAR
%   main.tex        -> preámbulo, datos de contacto y orden de las secciones
%   src/*.tex       -> contenido de cada sección, editable de forma independiente
%
% Para compilar:  xelatex main.tex
%==============================================================================

\documentclass[10pt,a4paper,sans]{moderncv}

% --- Estilo y color -----------------------------------------------------------
\moderncvstyle{banking}   % Estilo sobrio: sin barras laterales ni fotografía
\moderncvcolor{black}     % Monocromo: máxima legibilidad para filtros ATS

% --- Tipografía (requiere XeLaTeX/LuaLaTeX) -----------------------------------
\usepackage{fontspec}
% Las fuentes viajan dentro del repositorio (carpeta fonts/), así que el CV
% compila igual en Overleaf, en CI o en cualquier máquina, sin instalar nada.
% OJO: la opción «sans» de moderncv apunta \familydefault a \sfdefault, así que
% definir solo \setmainfont deja el cuerpo en Latin Modern. Hay que fijar ambas.
\defaultfontfeatures{
  Path            = fonts/,
  Scale           = 0.90,
  Extension       = .ttf,
  UprightFont     = *-Regular,
  BoldFont        = *-Bold,
  ItalicFont      = *-Italic,
  BoldItalicFont  = *-BoldItalic
}
\setmainfont{NotoSans}
\setsansfont{NotoSans}

\usepackage{polyglossia}
\setdefaultlanguage{spanish}

\usepackage[top=1.3cm,bottom=1.0cm,left=1.5cm,right=1.5cm]{geometry}
\setlength{\hintscolumnwidth}{2.2cm}

% --- Encabezado con mas aire -------------------------------------------------
% El estilo «banking» imprime nombre y cargo pegados en la misma linea, ambos en
% \Huge, separados por «|». Aqui se baja el cargo a su propia linea y se reduce
% su tamano, para que haya jerarquia visual entre nombre, cargo y contacto.
\makeatletter
\renewcommand*{\namefont}{\Huge\bfseries\upshape}
\renewcommand*{\titlefont}{\Large\mdseries\upshape}
% Cargo en su propia línea, mas pequeno que el nombre y con aire debajo.
\renewcommand*{\titlestyle}[1]{\\[0.7em]{\titlefont\textcolor{color2}{\@title}}\vspace{0.9em}}
% Datos de contacto con mas interlineado para que no se vean apelmazados.
\renewcommand*{\addressfont}{\fontsize{10}{16}\selectfont\mdseries\upshape}
% El bloque de contacto usa 0.8\textwidth y deja GitHub colgando en una linea
% suelta; ensancharlo al ancho completo lo equilibra. Si el parche falla, el
% encabezado sigue funcionando igual que antes.
\usepackage{etoolbox}
\patchcmd{\makehead}{0.8\textwidth}{\textwidth}{}{\typeout{AVISO: no se pudo ensanchar el bloque de contacto}}
\makeatother

% --- Interlineado compacto en las listas, para que el CV quepa en una página --
% Espaciado de secciones mas ajustado. moderncv reserva 2.5ex antes y 1ex
% despues de cada titulo de seccion; con seis secciones eso suma casi cuatro
% lineas. Esta es la definicion original de moderncvbodyiii, solo con los dos
% \addvspace reducidos.
\makeatletter
\RenewDocumentCommand{\section}{sm}{%
  \par\addvspace{1.6ex}%
  \phantomsection{}%
  \addcontentsline{toc}{section}{#2}%
  \strut\sectionstyle{#2}%
  \sectionrule%
  \par\nobreak\addvspace{0.7ex}\@afterheading}
\makeatother

\usepackage{needspace}   % evita que una sección quede partida entre páginas
\usepackage{enumitem}
\setlist[itemize]{topsep=1pt,itemsep=0pt,parsep=0pt,partopsep=0pt}

% --- Metadatos del PDF (los leen varios parsers de hojas de vida) -------------
% moderncv carga hyperref por su cuenta en \begin{document}; declararlo aquí con
% \usepackage provoca «Option clash for package hyperref». Por eso la
% configuración se aplica en un gancho que corre despues de esa carga.
\AtBeginDocument{%
  \hypersetup{
    pdftitle={Juan José Jiménez Guardo — Curriculum Vitae},
    pdfauthor={Juan José Jiménez Guardo},
    pdfsubject={Ingeniería de Sistemas — Análisis Funcional y Arquitectura de Software},
    pdfkeywords={Análisis Funcional, Arquitectura de Software, SAP, ERP, Laravel, Vue, Scrum, Modelado de Datos},
    hidelinks
  }%
}

%==============================================================================
% DATOS PERSONALES DE CONTACTO
%==============================================================================
\name{Juan José}{Jiménez Guardo}
\title{Arquitectura de Software \textbullet{} Desarrollador Full-Stack \textbullet{} Integración ERP/SAP}
\address{Cartagena de Indias}{Bolívar, Colombia}{}
\phone[mobile]{(+57) 300 8368682}
\email{juanjimenezg207@gmail.com}
\social[linkedin]{juan-jimenez-guardo}
\social[github]{JuanJimenezGuardo}

%==============================================================================
% DOCUMENTO
%==============================================================================
\begin{document}

% Sin numeración de páginas: moderncv activa el pie «1/2» en su propio
% \AtBeginDocument, de modo que hay que anularlo aquí, ya dentro del cuerpo.
\fancypagestyle{plain}{\fancyhf{}}
\pagestyle{empty}

\makecvtitle
\vspace*{-6mm}

% El orden de estos \input define el orden de las secciones del CV.
% Para reordenar o silenciar una sección, mueve o comenta su línea.
%==============================================================================
% Perfil Profesional
%==============================================================================
\section{Perfil Profesional}

Estudiante de último semestre de Ingeniería de Sistemas, orientado al análisis
funcional, la arquitectura de software y la integración de sistemas corporativos
(ERP/SAP). Actúo como puente entre operaciones industriales complejas —manufactura
ETO, entorno naval— y los equipos de desarrollo bajo metodologías ágiles (Scrum).
Mi aporte no se limita a la codificación: levanto reglas de negocio, diseño el
modelo relacional y defino la arquitectura de integración para que la solución
técnica resuelva el problema operativo real.

%==============================================================================
% Experiencia Profesional
%==============================================================================
\section{Experiencia Profesional}

\cventry{Agosto 2026--Presente}
        {Analista Funcional y Desarrollador de Software}
        {COTECMAR}
        {Cartagena, Colombia}{}
{Corporación de Ciencia y Tecnología para el Desarrollo de la Industria Naval.
\begin{itemize}
  \item Análisis funcional y diseño de la arquitectura de \textbf{TOP} (Portal de
        Operaciones Tecnológicas), sistema tipo MES que opera entre el piso de
        producción del astillero y el ERP corporativo SAP, en reemplazo de una
        gestión fragmentada en archivos Excel y cronogramas aislados.
  \item Ingeniería inversa del sistema \emph{legacy} y rediseño del modelo relacional
        (ERD documentado en Mermaid.js): unificación de estructuras redundantes de
        costos en una jerarquía única y corrección del modelo para soportar la
        realidad comercial de un proyecto con múltiples cascos.
  \item Diseño del modelo de integración con \textbf{SAP MM/PS}: inyección de
        solicitudes de pedido y reservas desde TOP, con una regla de caducidad que
        elimina las reservas fantasma responsables de inmovilizar inventario.
  \item Definición del indicador operativo de productividad (\textbf{Kg/HH}) como
        valor derivado, calculado sobre horas-hombre efectivamente auditadas y
        aprobadas, y del circuito de validación Supervisor $\rightarrow$ Jefe de Área
        previo a la imputación de consumos.
  \item Rediseño funcional del módulo comercial de estimaciones: máquina de ocho
        estados con la firma como prerrequisito del contrato, y vistas Kanban,
        tabla y calendario.
  \item Desarrollo sobre \textbf{Laravel 10} (PHP 8), \textbf{Vue 3} e
        \textbf{Inertia.js}, con MySQL como base operativa y conexiones SQL Server
        hacia SAP.
\end{itemize}}

\cventry{Agosto 2025--Mayo 2026}
        {Desarrollador Full-Stack --- Proyecto de Ingeniería de Sistemas}
        {VIT Cloud Services}
        {Bogotá, Colombia}{}
{\begin{itemize}
  \item Desarrollo full-stack de una plataforma web para la gestión de la
        Seguridad de la Información corporativa (SGSI), alineada a la norma
        \textbf{ISO/IEC 27001:2022}.
  \item Lideré la definición de requerimientos funcionales en comunicación
        directa con la gerencia del cliente, traduciendo los controles de la
        norma en flujos de trabajo y modelo de datos.
  \item Construcción con \textbf{Django}, \textbf{PostgreSQL} y \textbf{React}
        a lo largo de dos semestres académicos.
\end{itemize}}

% --- Duplicar este bloque para experiencias anteriores, prácticas o freelance ---
% \cventry{[AAAA]--[AAAA]}{[Cargo]}{[Empresa]}{[Ciudad]}{}{[Logros medibles]}

%==============================================================================
% Formación Académica
%==============================================================================
\section{Formación Académica}

\cventry{2022--Presente}
        {Ingeniería de Sistemas}
        {Universidad Tecnológica de Bolívar}
        {Cartagena, Colombia}{}
        {Énfasis en ingeniería de software, bases de datos y arquitectura de sistemas de información.}

\subsection{Certificaciones}

\cvlistitem{AWS Academy Graduate — Cloud Foundations.}
\cvlistitem{Bootcamp de Inteligencia Artificial, Nivel Intermedio — MinTIC / Universidad Tecnológica de Bolívar.}

% --- Para añadir otra certificación, duplica una línea \cvlistitem ---

%==============================================================================
% Competencias Técnicas
%==============================================================================
\section{Competencias Técnicas}

\cvitem{Arquitectura}{Modelado relacional y ERD, diccionarios de datos, diagramación con Mermaid.js, diseño transaccional, integridad referencial, arquitectura de integración entre sistemas.}
\cvitem{Backend}{PHP 8, Laravel 10, Eloquent ORM, diseño de APIs REST, migraciones versionadas, control de permisos y auditoría.}
\cvitem{Frontend}{Vue 3, Inertia.js, Vite, Tailwind CSS, JavaScript.}
\cvitem{Bases de datos}{MySQL, SQL Server, SQL avanzado, normalización, optimización de consultas.}
\cvitem{ERP / Integración}{SAP ERP (módulos MM y PS), homologación de maestros, orquestación de transacciones hacia el ERP, lógica de industria naval (ETO, grafos, SWBS).}
\cvitem{Metodologías}{Scrum, levantamiento de requerimientos funcionales, documentación técnica, gestión de proyectos TI.}
\cvitem{Herramientas}{Git, GitHub, Docker, Jenkins (CI/CD), \LaTeX{}.}

%==============================================================================
% Idiomas
%==============================================================================
\Needspace*{3\baselineskip}   % la sección no se parte entre páginas
\section{Idiomas}

\cvitem{Español}{Nativo}
\cvitem{Inglés}{B2}

% VIT paso a Experiencia Profesional, asi que esta seccion solo contendria TOP,
% que ya se describe en la entrada de COTECMAR. Se deja el archivo en el repo:
% para recuperarla basta descomentar la linea siguiente.
% %==============================================================================
% Proyectos Destacados
%==============================================================================
\section{Proyectos Destacados}

\cvitem{TOP}{Portal de Operaciones Tecnológicas de COTECMAR. Plataforma de ejecución de manufactura para un astillero mixto —construcción y reparación— que orquesta programación, planillación, productividad y el ciclo logístico de materiales, manteniendo SAP como fuente maestra de la verdad financiera y de inventario.}
\cvitem{VIT}{Plataforma SGSI alineada a la norma \textbf{ISO/IEC 27001:2022}. Desarrollo full-stack de una aplicación web para la gestión de la seguridad de la información corporativa. Lideré la definición de requerimientos y la construcción con \textbf{Django}, \textbf{PostgreSQL} y \textbf{React}, en comunicación directa con la gerencia del cliente (VIT Cloud Services).}

% \cvitem{[Proyecto]}{[Problema, solución y resultado en una o dos líneas.]}


\end{document}
